\documentclass[a4paper]{article}
\usepackage{amsfonts}
\usepackage{amsmath}
\usepackage{amssymb}
\usepackage{graphicx}     

\begin{document}
  
\noindent \large Auteur: Stan Van Campenhout \\
\noindent \large Studierichting:  Informatica\\
\noindent \large Oefeningenreeks 4A nr. 46.2\\

\medskip

\normalsize

$ \displaystyle\int\dfrac{-3x + 3}{2x^4+5x^2+2}dx = \displaystyle\int\dfrac{-3x+3}{\left(x^2+2\right)\left(2x^2+1\right)}dx
= \dfrac{1}{2}\displaystyle\int\dfrac{-3x+3}{\left(x^2+2\right)\left(x^2+\dfrac{1}{2}\right)}dx $\\

$ \rightarrow \dfrac{-3x+3}{\left(x^2+2\right)\left(x^2+\dfrac{1}{2}\right)} = \dfrac{Ax + B}{x^2+2} + \dfrac{Cx + D}{x^2 + \dfrac{1}{2}} $\\

$ A\left(\sqrt{2}i\right) + B = \dfrac{-3\sqrt{2}i + 3}{-2 + \dfrac{1}{2}} = -2 -2\sqrt{2}i $\\

$ \Rightarrow A = 2, B = -2 $\\

$ \Rightarrow C\left(\dfrac{\sqrt{2}}{2}i\right) + D = \dfrac{\dfrac{-3\sqrt{2}}{2}i + 3}{-\dfrac{1}{2}+2} = 2 - \sqrt{2}i $\\

$ \Rightarrow C = -2, D = 2 $\\

$ = \displaystyle\int\dfrac{x-1}{x^2+2}dx - \displaystyle\int\dfrac{x-1}{x^2+\dfrac{1}{2}}dx $\\

$ = \dfrac{1}{2}\displaystyle\int\dfrac{1}{x^2+2}d\left(x^2+2\right) - \displaystyle\int\dfrac{1}{x^2+2}dx -\dfrac{1}{2}\displaystyle\int\dfrac{1}{x^2+\dfrac{1}{2}}d\left(x^2+\dfrac{1}{2}\right) + \displaystyle\int\dfrac{1}{x^2+\dfrac{1}{2}}dx $\\

$ = \dfrac{1}{2}\ln|x^2+2|-\dfrac{1}{\sqrt{2}}\operatorname{Bgtan}\left(\dfrac{x}{\sqrt{2}}\right)-\dfrac{1}{2}\ln\left|x^2+\dfrac{1}{2}\right|+\sqrt{2}\operatorname{Bgtan}\left(\sqrt{2}x\right) + c $\\

\begin{thebibliography}{999}
%\bibitem{a} Referentie
\end{thebibliography}
\end{document} 
